% PAQ-KV — CROSS-READER GENERALITY companion note (internal; NOT the submission).
% Standalone document, separate from paq_aaai.tex. Consolidates the multi-VLM
% second-reader study for internal review; the team decides later what (if any)
% folds into the main manuscript.
% Every number is taken verbatim from the verified run artifacts (see App. A).
% Build: ./tectonic -X compile paq_cross_reader_generality.tex
% ============================================================================
\documentclass[letterpaper]{article} % DO NOT CHANGE THIS
% ---------------------------------------------------------------------------
% LOCAL BUILD SHIM (mirrors paq_aaai.tex): host has no pdflatex; build with
% tectonic (XeTeX). Remove the shim for any official pdflatex build.
\usepackage{iftex}
\ifPDFTeX\else\let\RequirePDFTeX\relax\fi
% ---------------------------------------------------------------------------
\usepackage[submission]{aaai2027}  % DO NOT CHANGE THIS
\usepackage[hyphens]{url}  % DO NOT CHANGE THIS
\usepackage{graphicx} % DO NOT CHANGE THIS
\urlstyle{rm} % DO NOT CHANGE THIS
\def\UrlFont{\rm}  % DO NOT CHANGE THIS
\usepackage{natbib}  % DO NOT CHANGE THIS AND DO NOT ADD ANY OPTIONS TO IT
\usepackage{caption} % DO NOT CHANGE THIS AND DO NOT ADD ANY OPTIONS TO IT
\frenchspacing  % DO NOT CHANGE THIS
\usepackage{booktabs}
\usepackage{multirow}
\usepackage{amsmath}
\usepackage{amssymb}

\pdfinfo{
/TemplateVersion (2027.1)
}

\setcounter{secnumdepth}{2}

\newcommand{\paq}{\textsc{PAQ}}
\newcommand{\paqc}{\textsc{PAQ-KV}}
\newcommand{\lb}{\mathrm{LB}_{95}}
\newcommand{\pending}{\textit{pending}}

\title{Cross-Reader Generality of \paqc{}:\\A Multi-VLM Study (Internal Companion Note)}
\author{Anonymous submission}
\affiliations{Anonymous submission}

\begin{document}
\maketitle

\begin{abstract}
\noindent
\paqc{} is a training-free, index-free, single-reader method: a vision--language model (VLM) prefills a multi-page
document once, and per query computes its \emph{own} question-to-context attention, selects the top $\approx$5\% of
64-token visual key/value (KV) blocks, and decodes under a reversible 4-D restricted-attention mask over its own KV
cache. Its headline result---beating a trained ColQwen2 retriever on MMLongBench-Doc accuracy---is established on
Qwen3-VL-8B and confirmed out-of-sample. This note asks the natural follow-up: \emph{does the advantage generalize to
other VLM readers?} We port the full method (reader-own scoring, selection, and masked decode) to additional VLMs and
report matched-budget results on a fixed 80-document / 444-query development (DEV) pin. Findings so far: on
Qwen3-VL-4B (same architecture, half the parameters) the method largely \emph{holds} at byte-identical budget
($\Delta=+0.030$, $\lb=-0.004$, a near-win that denoises above full context); on Qwen2.5-VL-7B (same Qwen-VL family) it
is a \emph{statistical dead heat} ($\Delta=+0.007$, $\lb=-0.037$); on InternVL3.5-8B (a different VLM whose language
backbone is Qwen3) it \emph{significantly underperforms} ($\Delta=-0.067$, $\lb=-0.106$). Together these characterize a
\textbf{mechanistic boundary} located at architecture rather than scale: the method's benefit tracks how good the host
reader's own attention is as a below-page retriever. Three further readers (Molmo-7B-D, MiniCPM-V-4.5, PaliGemma-2)
were each found incompatible with the current setup for distinct, honest reasons (remote-code version drift for the
first two; an $\approx$8K context that cannot hold a multi-page document for the third), so the study's evidence rests
on the four readers above. This is an internal companion note kept separate from the submission; the shared team
decides what, if anything, to fold in.
\end{abstract}

\section{Setup}
\label{sec:setup}

\subsection{The method and the established result}
\paqc{} prefills a full multi-page document once into a VLM's KV cache. For each query it (i) reads the reader's own
question-to-context attention (scorer \textsc{v0}: mean over heads and layers), (ii) partitions the retained visual KV
into 64-token stripe blocks and selects the top $\approx$5\% by attention score, and (iii) decodes the answer under a
reversible 4-D additive attention mask that hides all non-selected visual columns while always keeping the (small)
non-visual scaffold. It trains nothing, builds no external index, and uses a single model.

On \textbf{Qwen3-VL-8B}, out-of-sample on a sealed 45-document / 244-query MMLongBench-Doc split, \paqc{} at 5\%
scores \textbf{0.3957} versus a trained ColQwen2 retriever's \textbf{0.3202} ($+0.0755$; doc-cluster $\lb=+0.0290$).
On the 80-document / 444-query DEV pin the same beat is $+0.079$ ($\lb=+0.041$). This note does not touch that result;
the sealed split remains evaluated exactly once.

\subsection{What generality requires of a reader}
\label{sec:requires}
Porting surfaced hard prerequisites that any candidate reader must satisfy for the method to even run:
\begin{enumerate}
\item \textbf{White-box access.} The method reads attention weights, edits the KV cache, and injects a 4-D mask. Any
API / closed-weights model is categorically impossible.
\item \textbf{In-line visual tokens in the decoder's self-attention KV.} Blocks are partitions of visual token
\emph{positions in the self-attention sequence}. A reader that injects vision via \emph{cross-attention} (vision lives
outside the self-attention KV) has no visual KV blocks to select or mask.
\item \textbf{Hookable eager attention returning 4-D weights}, plus a \textbf{4-D additive mask honored by the
language model} (the transferability crux; native \textsc{transformers} classes that delegate to a native decoder
satisfy this, custom \texttt{trust\_remote\_code} attentions may not).
\item \textbf{Reconstructable, continuable positions} (1-D \texttt{arange} or M-RoPE) and enough context to prefill a
multi-page document. Sub-native context ceilings force a reduced-budget setting (a disclosed confound,
\S\ref{sec:confounds}).
\item Empirically (\S\ref{sec:discussion}), \textbf{the reader's own attention must actually be a good below-page
retriever.} Mechanical portability is necessary but not sufficient.
\end{enumerate}

\subsection{Protocol}
\label{sec:protocol}
All second-reader numbers are on the fixed \textbf{80-document / 444-query DEV pin} (\texttt{paq\_confirm\_pin.json},
MMLongBench-Doc DEV), reported as a \emph{generality control}, not a new sealed claim. Each reader is compared to
itself under two selectors at a matched budget: \paqc{} (own attention selects $\approx$5\% of blocks) vs.\ a
\emph{comparator} arm in which the same reader reads the top-5\% pages chosen by cached ColQwen2 retrieval scores
(no new ColQwen2 runs). We also record a \emph{full-context} reference. The primary statistic is the doc-cluster
95\% lower bound $\lb$ of $(\textrm{paqkv}-\textrm{comparator})$ (seed 0, 10{,}000 bootstrap resamples); the
pre-registered ``transfers'' bar is $\lb>0$. An adversarial worst-case imputation of any missing query binds the
verdict. A hard \textbf{identity gate} (teacher-forced argmax parity of the all-true 4-D mask vs.\ native causal
decode, plus a restrictive mask changing the output, on $\geq$2 distinct documents including one $\geq$30 pages) must
pass before any number is trusted. Verdicts are labeled \textsc{incompatible} (identity fails / port infeasible),
\textsc{incomplete} (coverage/power short), \textsc{reader2\_no\_transfer} ($\lb\leq 0$ with clean coverage), or
\textsc{reader2\_transfers} ($\lb>0$).

\section{Results}
\label{sec:results}

\subsection{Main comparison (verified readers)}
Table~\ref{tab:main} reports the three verified readers on the same 444-query DEV pin. Accuracy is the MMLongBench-Doc
scorer; higher is better. All three readers read the same content selectors (own \paqc{} blocks vs.\ ColQwen2 pages)
at a matched $\approx$5\% budget.

\begin{table*}[t]
\centering
\caption{\paqc{} vs.\ the ColQwen2-page comparator, within each reader, on the 80-doc / 444-q DEV pin.
$\Delta=$ mean(paqkv $-$ ColQwen2); $\lb=$ doc-cluster 95\% lower bound (pre-registered ``transfers'' bar is
$\lb>0$). ``ctx'' is the mean realized prefill length. Qwen3-VL-8B is the reader the method was developed and
sealed on; Qwen3-VL-4B uses byte-identical input geometry (a pure size ablation).}
\label{tab:main}
\small
\begin{tabular}{@{}lcccccc@{}}
\toprule
Reader & ctx & \paqc{} & ColQwen2 & Full & $\Delta$ & $\lb$ \\
\midrule
Qwen3-VL-8B \textit{(native, sealed)}   & $\sim$69K & \textbf{0.4042} & 0.3284 & 0.3965 & $+0.076$ & $+0.041$ \\
Qwen3-VL-4B \textit{(size ablation)}    & $\sim$69K & 0.3640 & 0.3336 & 0.3480 & $+0.030$ & $-0.004$ \\
Qwen2.5-VL-7B                           & $\sim$52K & 0.2698 & 0.2623 & 0.3237 & $+0.007$ & $-0.037$ \\
InternVL3.5-8B                          & $\sim$36K & 0.1876 & 0.2542 & 0.2637 & $-0.067$ & $-0.106$ \\
\bottomrule
\end{tabular}
\end{table*}

\begin{itemize}
\item \textbf{Qwen3-VL-8B (native reader): win.} $\paqc{}-\textrm{ColQwen2}=+0.076$, $\lb=+0.041$; \paqc{} even edges
its own full context ($0.4042$ vs.\ $0.3965$), i.e.\ selection \emph{denoises}.
\item \textbf{Qwen3-VL-4B (same architecture, half the parameters): near-win / method holds.} $\Delta=+0.030$,
$\lb=-0.004$ (just misses the bar at $n=444$; the point estimate is a clear positive). \paqc{} again \emph{denoises}
above its own full context ($0.3640$ vs.\ $0.3480$), and on the 137 answerable queries it recovers \textbf{0.869} vs.\
the comparator's \textbf{0.720}---the strongest below-page selection of any second reader. Per-query: 63 wins / 336
ties / 45 losses. Because the input geometry is byte-identical to the 8B (verified on all 80 documents), the gap to
the 8B ($-0.040$ cross-reader) is a pure capacity effect, not a budget artifact. The method's mechanism clearly
survives shrinking the reader within the same architecture.
\item \textbf{Qwen2.5-VL-7B (same Qwen-VL family): tie.} $\Delta=+0.0074$, $\lb=[-0.037,+0.049]$---a statistical dead
heat straddling zero (\paqc{} nominally ahead). Verdict \textsc{reader2\_no\_transfer} by the $\lb>0$ bar, but this is
``no detectable difference,'' not a loss. Per-query: 54 wins / 331 ties / 59 losses. On the 128 queries answerable by
full context, \paqc{} recovers \textbf{0.703} vs.\ the comparator's \textbf{0.581}---here \paqc{} is the \emph{better}
below-page selector. It loses only $-0.054$ $[-0.087,-0.024]$ to its own full context.
\item \textbf{InternVL3.5-8B (cross-family, Qwen3 LM): significant negative.}
$\paqc{}-\textrm{ColQwen2}=-0.067$, $\lb=[-0.106,-0.029]$ (whole interval below zero). It also trails its own full
context ($-0.076$ $[-0.113,-0.041]$). On answerable queries it recovers only \textbf{0.529}, losing $\approx$47\% of
them to its 8\% selection. Verdict \textsc{reader2\_no\_transfer} (a genuine, well-powered negative).
\end{itemize}

\subsection{Additional readers (in progress)}
\label{sec:additional}
Four further readers were requested for this study. Qwen3-VL-4B is \textbf{complete} and appears in
Tables~\ref{tab:main} and \ref{tab:gradient} (it is the same \textsc{transformers} class as the sealed 8B, so it
reused the audited harness with a loader swap and reproduced the 8B's byte-identical input geometry). The remaining
three (Table~\ref{tab:additional}) each require a genuine per-model port because their attention APIs differ from the
standard \textsc{transformers} 4-D-mask path the method relies on. Feasibility is decided cheaply by a Phase-0 identity
spike before any full run; full ports proceed only for spike-passers, and a reader that fails the identity gate is
reported honestly as \textsc{incompatible} rather than forced.

\begin{table*}[t]
\centering
\caption{Additional candidate readers. ``Native'' = a \textsc{transformers} 5.5 class exists (mask machinery and eager
hooks transfer as for InternVL/Qwen2.5-VL); ``remote'' = \texttt{trust\_remote\_code} custom attention (4-D mask honored
and hookability must be validated by the spike). Results are filled in as runs complete.}
\label{tab:additional}
\small
\setlength{\tabcolsep}{5pt}
\begin{tabular}{@{}lllcc@{}}
\toprule
Reader & LM backbone & Vision integration & Impl. & Phase-0 verdict \\
\midrule
Molmo-7B-D    & Qwen2-7B & CLIP, in-line tokens  & remote & \textsc{incompatible} (model won't load: version drift) \\
MiniCPM-V-4.5 & Qwen3-8B & 3D-Resampler (64 tok) & remote & \textsc{incompatible} (processor won't build inputs) \\
PaliGemma-2   & Gemma-2  & prefix-LM (bidir.)    & native & \textsc{incompatible} ($\sim$8K ctx: cannot fit a doc) \\
\bottomrule
\end{tabular}
\end{table*}

\paragraph{Per-reader notes.}
\begin{itemize}
\item \textbf{Qwen3-VL-4B} is the same architecture as the sealed 8B (same \textsc{transformers} class, M-RoPE, native
per-image budget), so the existing harness runs with a loader swap. It is a \emph{within-architecture size ablation},
not a cross-architecture test; a win here would say the method survives shrinking the same reader, not that it
transfers across families.
\item \textbf{Molmo-7B-D} $\to$ \textsc{incompatible} (version drift, not mechanism). Its custom attention
\emph{does} expose the exact injection point the method needs---the top-level \texttt{forward} accepts a 4-D
\texttt{attention\_bias} fed straight to SDPA---so the mechanism is portable in principle. But its 2024 remote-code
modeling file will not load under the pinned \textsc{transformers} 5.5.4: a cascade of tied-weights API mismatches
(\texttt{all\_tied\_weights\_keys}, then \texttt{tie\_weights(missing\_keys=...)}) blocks \texttt{from\_pretrained}.
Loading it would require porting Molmo's modeling file to the current API, out of scope here. This is an
engineering/version-compatibility incompatibility, not a statement about the method.
\item \textbf{MiniCPM-V-4.5} $\to$ \textsc{incompatible} (version drift, not mechanism). The model weights load, but
its 2025 remote-code \emph{processor} cannot build inputs under the pinned \textsc{transformers} 5.5.4: it expects a
custom tokenizer attribute (\texttt{tokenizer.im\_start\_id}) that the current fast-tokenizer backend does not expose.
Independently, its 3D-Resampler compresses each page to $\approx$64 tokens, so even a working port would be
near-degenerate below page granularity (a 5\% block budget of $\approx$64 tokens is $<1$ block per page); and its LM
is Qwen3, the same family as the InternVL negative.
\item \textbf{PaliGemma-2} $\to$ \textsc{incompatible} (representation, not version). It loads natively and accepts
multiple images (1024 tokens/page), but its Gemma-2 language model has an $\approx$8192-token context with a 4096
sliding window, and PaliGemma uses a \emph{fixed} 1024 tokens per image with no dynamic per-page budget. That caps a
document at $\approx$7 pages; the study's documents are 17--104 pages (median $\approx$32), so \paqc{}'s ``prefill the
whole document once'' premise cannot be satisfied---the document does not fit. (Its prefix-LM attention and
short-output tuning would further complicate a port even at fitting lengths.)
\end{itemize}

\noindent\textbf{Takeaway from the three incompatibilities.} They are not method failures; they map the method's
\emph{practical} deployment envelope. Running reader-own \paqc{} needs, in addition to the mechanistic property of
\S\ref{sec:discussion}: (i) a version-current implementation exposing a hookable attention and an honored additive
mask; and (ii) enough context, with a dynamic per-page token budget, to prefill a whole multi-page document. Among the
readers surveyed, the Qwen-VL family (Qwen3-VL-8B/4B, Qwen2.5-VL-7B) and InternVL3.5-8B satisfy both; the three above
each miss one. This is worth stating plainly as a scope condition rather than omitting it.

\section{Discussion: a mechanistic generality gradient}
\label{sec:discussion}
The verified readers order cleanly by architectural distance from the reader the method was developed on
(Table~\ref{tab:gradient}). The benefit is not universal; it \emph{degrades monotonically} as the reader moves away
from Qwen3-VL-8B, and the degradation tracks a measurable property---how well the reader's own question-to-context
attention localizes the answer below page granularity. Qwen3-VL-8B has this property (it recovers answerable queries
well at 5\% and even denoises above full context); \textbf{Qwen3-VL-4B retains it} (near-win at matched, byte-identical
budget; denoises above full; recovers 0.87 of answerable queries), so the effect survives halving the parameter count
\emph{within the same architecture} and is not an artifact of the 8B's scale alone; Qwen2.5-VL-7B partially has it
(break-even; recovers 0.70); InternVL3.5-8B lacks it here (recovers 0.53 and loses). The two Qwen3-VL points (win,
near-win) versus the same-family tie and the cross-family negative locate the boundary at \emph{architecture}, not
merely scale. This \emph{confirms} rather than refutes the single-reader observation, and reframes it from an
unexamined weakness into an investigated boundary condition: \paqc{} is a white-box method whose benefit is contingent
on the host reader's attention being an effective below-page retriever.

\begin{table*}[t]
\centering
\caption{The generality gradient, ordered by architectural distance from the sealed reader.
$\Delta$ / doc-cluster $\lb$ (paqkv $-$ comparator), DEV pin (Qwen3-VL-8B also shows a sealed out-of-sample
$\lb=+0.029$). The benefit degrades monotonically with distance.}
\label{tab:gradient}
\small
\begin{tabular}{@{}llc@{}}
\toprule
Reader & Relation to the sealed reader & $\Delta$ / $\lb$ \\
\midrule
Qwen3-VL-8B    & native (sealed reader)      & $+0.076$ / $+0.041$ \textit{(win)} \\
Qwen3-VL-4B    & same architecture, 4B       & $+0.030$ / $-0.004$ \textit{(near-win)} \\
Qwen2.5-VL-7B  & same family, prior gen.\    & $+0.007$ / $-0.037$ \textit{(tie)} \\
InternVL3.5-8B & cross-family VLM (Qwen3 LM) & $-0.067$ / $-0.106$ \textit{(neg.)} \\
\bottomrule
\end{tabular}
\end{table*}

\section{Confounds and honesty caveats}
\label{sec:confounds}
Any use of these numbers must carry the following.
\begin{enumerate}
\item \textbf{Reduced-budget confound.} Both second readers run below the native-long 8B budget (InternVL $\sim$36K via
a 40K context ceiling; Qwen2.5-VL $\sim$52K via a 128K ceiling plus a processor pre-resize). Within-reader
matched-budget comparisons carry the findings; cross-reader comparisons \emph{against} the 8B are confounded by
capacity and budget and are secondary only.
\item \textbf{Shared LM family.} InternVL3.5-8B's language backbone is Qwen3, which is \emph{why} the KV/attention/mask
machinery transfers. It is a legitimate cross-VLM result but not a fully independent language architecture; disclose it.
\item \textbf{Conservative-against-method budget asymmetry (Qwen2.5-VL).} The method downscales all pages in the
full-document prefill, whereas the comparator renders its 1--5 kept pages near-native and thus gets a slightly larger
per-page budget. The tie is therefore a floor on the within-family effect.
\item \textbf{DEV-pin, not a new sealed claim.} All second-reader numbers are on the DEV pin as a generality control.
The sealed split is untouched.
\end{enumerate}

\section*{Verification (Appendix A)}
\label{sec:verification}
Every reported run passed a hard identity gate and was independently checked. For each verified reader: the identity
audit confirmed the all-true 4-D mask reproduces native causal decoding under the teacher-forced argmax criterion
(max $|\Delta\text{logit}|=0$), a restrictive mask changes the output, on $\geq$2 distinct documents (one $\geq$30
pages), with the resolved model revision and code hash bound into the frozen run manifest. A CPU sparse-mask audit
confirmed the block-to-column mapping is exact (all-blocks selection reproduces the all-true mask; a sparse subset
masks exactly the intended columns). An off-meter code reviewer independently attempted to falsify each result and
found no implementation defect (``genuine negative'' for InternVL; ``genuine null with favorable caveats'' for
Qwen2.5-VL). Full coverage for both second readers: 444/444 queries, 80 documents, 0 skipped, 0 dropped, identity pass.
Artifacts, verdict JSONs, reviewer transcripts, and audit scripts are in the project repository under
\texttt{research/} and \texttt{research/reviews/paq\_second\_vlm\_2026-07-22/}; the consolidated running log is
\texttt{research/design\_notes/paq\_cross\_reader\_generality\_2026-07-23.md}.

\section*{References (model / benchmark sources)}
\begin{description}
\item[Qwen3-VL / Qwen2.5-VL] Qwen Team. \url{https://huggingface.co/Qwen/Qwen2.5-VL-7B-Instruct}
\item[InternVL3.5-8B] OpenGVLab. \url{https://huggingface.co/OpenGVLab/InternVL3_5-8B-HF}
\item[Molmo-7B-D] Allen Institute for AI. \url{https://huggingface.co/allenai/Molmo-7B-D-0924}
\item[MiniCPM-V-4.5] OpenBMB. \url{https://huggingface.co/openbmb/MiniCPM-V-4_5}
\item[PaliGemma-2] Google. \url{https://huggingface.co/google/paligemma2-3b-mix-448}
\item[ColQwen2] Faysse et al. \url{https://huggingface.co/vidore/colqwen2-v1.0}
\item[MMLongBench-Doc] Ma et al. \url{https://mmlongbench-doc.github.io}
\end{description}

\end{document}
